\documentclass[11pt,a4paper]{article}

\usepackage[utf8]{inputenc}
\usepackage[T1]{fontenc}
\usepackage{lmodern}
\usepackage[french]{babel}
\usepackage[margin=2.5cm]{geometry}
\usepackage{xcolor}
\usepackage{enumitem}
\usepackage{booktabs}
\usepackage{array}
\usepackage{titlesec}
\usepackage{fancyhdr}
\usepackage[hidelinks]{hyperref}

\definecolor{linkblue}{HTML}{1F5FBF}
\definecolor{linkgrey}{HTML}{555555}

\titleformat{\section}{\Large\bfseries\color{linkblue}}{\thesection}{0.6em}{}
\titleformat{\subsection}{\large\bfseries\color{linkblue}}{\thesubsection}{0.6em}{}

\pagestyle{fancy}
\fancyhf{}
\fancyhead[L]{\small\color{linkgrey} Link --- Présentation de l'application}
\fancyfoot[C]{\small\thepage}
\renewcommand{\headrulewidth}{0.3pt}

\setlength{\parindent}{0pt}
\setlength{\parskip}{0.5em}

% Une fonctionnalité : titre en gras + description courte
\newcommand{\fonc}[2]{\textbf{#1.} #2\par\vspace{0.2em}}

\begin{document}

	%-------------------------------------------------
	% PAGE DE TITRE
	%-------------------------------------------------
	\begin{titlepage}
		\centering
		\vspace*{3cm}
		{\fontsize{48}{56}\selectfont\bfseries\color{linkblue} Link}\\[0.8cm]
		{\Large Application mobile de suivi d'activité\\pour très petites entreprises}\\[5cm]
		{\large Document de présentation et de description fonctionnelle}\\[0.5cm]
		{\color{linkgrey} Version 1.2 --- Novembre 2026}
		\vfill
		{\color{linkgrey}\small NDEFFEU Arthur}
	\end{titlepage}

	\tableofcontents
	\newpage

	%-------------------------------------------------
	\section{Introduction}
	%-------------------------------------------------
	\textbf{Link} est une application mobile destinée aux très petites entreprises (cinq personnes au maximum) : salon de coiffure, boutique de vêtements, atelier, petit commerce de proximité, etc. Elle permet à un responsable (le \emph{boss}) et à ses employés d'enregistrer chaque vente ou service rendu, puis d'obtenir une vue claire de l'activité et un bénéfice indicatif.

	Link n'est pas un logiciel de gestion complet. Il n'y a ni comptabilité réglementaire, ni facturation, ni gestion de paie. L'objectif est un \textbf{suivi simple, rapide et flexible} de ce que l'entreprise produit, de ce que chacun y contribue et de ce qu'elle dépense.

	L'activité est organisée en \textbf{exercices} : le boss ouvre un exercice, y enregistre l'activité, puis le clôture manuellement pour en ouvrir un nouveau. Chez le boss, aucune opération n'est possible sans exercice ouvert.

	Sa particularité est d'être \textbf{pensée d'abord pour le hors ligne}. Chaque téléphone enregistre l'activité sans connexion. Les appareils se synchronisent ensuite avec celui du boss, à la demande de l'utilisateur, directement sur le réseau local (Wi-Fi de l'entreprise ou partage de connexion du téléphone du boss), sans besoin d'internet. Un serveur relais sur internet existe en complément, pour les cas où le boss est absent, mais l'application fonctionne entièrement sans lui.

	%-------------------------------------------------
	\section{Contexte et objectifs}
	%-------------------------------------------------
	\subsection{Contexte}
	Dans un très petit commerce, l'activité est souvent notée sur papier ou de mémoire. Les outils existants sont trop lourds, exigent une connexion permanente ou ne sont pas adaptés aux prix négociés et aux usages informels. Malgré l'omniprésence des téléphones portables, le problème de fiabilité de la connexion internet reste présent.

	\subsection{Objectifs}
	\begin{itemize}[leftmargin=*]
		\item Permettre d'enregistrer une vente ou un service en quelques secondes, même sans connexion, et d'en enregistrer plusieurs très rapidement en cas d'affluence (mode rafale).
		\item Donner au boss une vision globale et à jour de son entreprise, avec le chiffre d'affaires, les dépenses et un bénéfice indicatif.
		\item Permettre à chaque employé de suivre sa propre activité, si le boss le souhaite.
		\item Garder les données cohérentes entre tous les appareils grâce à la synchronisation, sans dépendre d'internet.
		\item Rester très flexible : prix libres, stock optionnel, modifications possibles à tout moment par le boss.
		\item Organiser l'activité en exercices clairement séparés.
		\item Protéger les données (chiffrement, sauvegarde, copie de secours, révocation d'appareil).
	\end{itemize}

	\subsection{Hors périmètre}
	Comptabilité réglementaire, TVA et facturation légale, reçus et factures, gestion de la paie, achats et fournisseurs, gestion détaillée du stock (lots, coût moyen pondéré, inventaire complet de tout le magasin en une seule opération). Les éléments volontairement mis de côté pour cette version sont listés à la section~\ref{sec:miscote}.

	%-------------------------------------------------
	\section{Acteurs et droits}
	%-------------------------------------------------
	Link distingue deux rôles. Le modèle de données prévoit dès le départ l'ajout d'autres rôles, mais ils ne font pas partie de cette version. Le boss participe lui aussi à la production : il peut vendre comme un employé.

	Un principe guide les droits de l'employé : \textbf{il ne voit jamais que ce qu'il a fait lui-même}, et seulement dans la mesure permise par le réglage de conservation choisi par le boss.

	\begin{center}
		\small
		\begin{tabular}{p{7.2cm} c c}
			\toprule
			\textbf{Action} & \textbf{Boss} & \textbf{Employé} \\
			\midrule
			Créer et configurer l'entreprise & oui & non \\
			Créer, désactiver, réinitialiser des comptes & oui & non \\
			Gérer le catalogue (produits, services, prix de vente) & oui & non \\
			Voir et modifier les prix d'achat & oui & non \\
			Voir les quantités en stock & oui & oui \\
			Modifier le stock (entrée, perte, inventaire) & oui & non \\
			Recevoir les alertes de stock & oui & non \\
			Créer et clôturer les exercices & oui & non \\
			Enregistrer une vente ou un service & oui & oui \\
			Utiliser le mode rafale & oui & oui \\
			Enregistrer une vente au nom d'un employé & oui & non \\
			Enregistrer et gérer les dépenses & oui & non \\
			Voir l'activité de toute l'équipe & oui & non \\
			Voir ses opérations en attente de synchronisation & oui & oui \\
			Revoir ses propres opérations déjà synchronisées, par exercice & oui & selon le réglage \\
			Consulter le contenu complet des exercices clos & oui & non \\
			Tableau de bord global et bénéfice & oui & non \\
			Choisir de conserver ou non l'historique des employés & oui & non \\
			Lancer une synchronisation & oui & oui \\
			Modifier son profil et son mot de passe & oui & oui \\
			Corriger une opération avant synchronisation & oui & oui (les siennes) \\
			Corriger ou annuler une opération après synchronisation & oui & non \\
			Révoquer un appareil & oui & non \\
			Gérer les sauvegardes & oui & non \\
			\bottomrule
		\end{tabular}
	\end{center}

	%-------------------------------------------------
	\section{Fonctionnalités détaillées}
	%-------------------------------------------------

	\subsection{Entreprise et comptes}

	\fonc{Création de l'entreprise}{Au premier lancement, le boss crée son compte puis son entreprise : nom, logo, couleurs, type d'activité et coordonnées. La devise est le franc CFA (FCFA) ; les montants sont des nombres entiers, sans virgule. Ces informations sont modifiables à tout moment, sauf la devise. Le boss crée ensuite son premier exercice (voir la section Exercices).}

	\fonc{Comptes employés}{Le boss crée les comptes (nom, prénom, identifiant, mot de passe provisoire). Il peut les modifier, les désactiver ou les réactiver. Un compte n'est jamais supprimé, afin de conserver l'historique. La création des comptes, comme celle des produits et des services, ne dépend pas de l'exercice.}

	\fonc{Première connexion}{Un employé doit remplacer son mot de passe provisoire dès sa première connexion. L'accès à l'application est bloqué tant que ce changement n'est pas fait.}

	\fonc{Réinitialisation de compte}{Si un employé oublie son mot de passe, le boss génère un nouveau mot de passe provisoire. Il est transmis à l'employé, qui doit le changer à la connexion suivante.}

	\fonc{Profil personnel}{Chaque utilisateur modifie ses informations : nom, prénom, téléphone, photo, langue et mot de passe. Les changements sont transmis au boss à la synchronisation suivante, puis aux autres appareils de l'utilisateur à leur propre synchronisation. D'ici là, un autre appareil de l'utilisateur garde l'ancien mot de passe. Si le profil a été modifié sur deux appareils, la dernière modification reçue par le boss l'emporte.}

	\subsection{Appareils et appairage}

	\fonc{Un téléphone, un utilisateur}{Chaque appareil appartient à un seul utilisateur. Un téléphone ne peut pas être partagé entre plusieurs employés : seul son propriétaire peut s'y connecter. Les opérations saisies sur un appareil sont donc toujours celles de son propriétaire.}

	\fonc{Appairage par QR code}{Pour rejoindre l'entreprise, un appareil doit être appairé. Le boss choisit le compte de l'employé concerné et son téléphone affiche un QR code, que le nouvel appareil scanne. Ce scan est obligatoire avant toute connexion, car le compte de l'employé est créé chez le boss. L'appareil est alors lié à ce compte et à lui seul. À cette occasion, les deux téléphones échangent leurs clés de chiffrement : l'appareil de l'employé reçoit notamment la \emph{clé publique} du boss, qui sert à la copie de secours (voir plus loin). L'appareil reçoit ensuite une copie des données qui le concernent : catalogue, paramètres et le compte de son propriétaire uniquement, jamais ceux des autres employés.}

	\fonc{Plusieurs appareils pour un même utilisateur}{Un employé peut avoir plusieurs appareils, chacun appairé séparément à son compte. Les appareils ne communiquent jamais entre eux : chacun envoie ses propres opérations au boss lors de sa propre synchronisation. Il revient à l'employé de garder chacun de ses appareils à jour. En mode conservation, chaque appareil ne garde que l'historique des opérations saisies sur lui, plus celles que le boss a saisies au nom de l'employé.}

	\fonc{Téléphone du boss}{Le téléphone du boss est le centre de l'entreprise. Le compte boss ne peut être utilisé que sur ce téléphone ; il ne peut pas se connecter sur l'appareil d'un employé.}

	\fonc{Révocation d'appareil}{Le boss peut révoquer un appareil (téléphone perdu, changé ou rendu). L'appareil révoqué efface toutes ses données locales, copie de secours comprise, à sa prochaine synchronisation. Désactiver un employé révoque tous ses appareils. Pour qu'un téléphone change de propriétaire, il est révoqué puis ré-appairé avec le nouveau compte.}

	\fonc{Perte d'un téléphone employé}{L'appareil est révoqué et un nouvel appareil est appairé. Aucun historique n'est restitué sur le nouvel appareil : l'employé repart avec les seules données reçues après le nouvel appairage.}

	\subsection{Exercices}

	\fonc{Notion d'exercice}{L'activité est découpée en exercices, définis librement par le boss. Un exercice a un nom et, si le boss le souhaite, une description. Il n'y a jamais qu'un seul exercice ouvert à la fois.}

	\fonc{Création du premier exercice}{Rien n'est créé automatiquement : le boss doit créer lui-même son premier exercice, une étape guidée lui rappelant que cette fonctionnalité existe. Tant qu'aucun exercice n'est ouvert, le boss ne peut enregistrer ni vente ni dépense, et son téléphone n'accepte aucune synchronisation. Les autres créations (comptes, produits, services, catégories) restent possibles.}

	\fonc{Clôture et nouvel exercice}{Le boss clôture l'exercice manuellement. L'action habituelle est « Clôturer et ouvrir un nouvel exercice », qui enchaîne les deux en donnant le nom (et éventuellement la description) du suivant. Entre une clôture et l'ouverture suivante, aucune synchronisation n'est possible. Avant de clôturer, l'application affiche, pour chaque employé, la date de sa dernière synchronisation et rappelle que les opérations qui n'ont pas encore été reçues compteront dans le prochain exercice.}

	\fonc{Rattachement d'une opération}{Une opération appartient à l'exercice \textbf{ouvert chez le boss au moment où il la reçoit}. Une opération saisie par le boss lui-même, ou une dépense, appartient à l'exercice ouvert au moment de sa saisie. Ainsi, une vente faite hors ligne avant une clôture mais synchronisée après rejoint le nouvel exercice. L'heure de l'opération est conservée pour l'historique, les statistiques par période et le stock, mais elle ne détermine pas l'exercice. Sur le téléphone de l'employé, une opération n'a pas encore d'exercice tant qu'elle n'a pas été reçue par le boss.}

	\fonc{Données concernées}{Les ventes et les dépenses appartiennent à un exercice. Le catalogue, les comptes, les clients, les paramètres et les quantités en stock sont communs à tous les exercices et ne sont pas remis à zéro. Les chiffres d'un exercice (chiffre d'affaires, dépenses, bénéfice) ne se reportent pas sur le suivant.}

	\fonc{Exercices clos}{Un exercice clos est strictement en lecture seule : il ne reçoit plus aucune opération et n'est plus pris en compte dans le tableau de bord de l'exercice courant. Il est archivé et n'est jamais supprimé. Le boss en consulte tout le contenu ; un employé, en mode conservation uniquement, n'y retrouve que ses propres opérations.}

	\subsection{Catalogue et stock}

	\fonc{Produits}{Le boss enregistre ses produits : nom, catégorie, prix de vente de référence, prix d'achat (facultatif), photo, description, référence interne. Un produit peut être archivé : il disparaît de la saisie mais reste dans l'historique.}

	\fonc{Prix d'achat}{Le prix d'achat est un champ facultatif du produit, indépendant du stock. Il sert à estimer le coût des articles vendus. Il n'est visible que du boss et n'est jamais envoyé aux employés. Un produit sans prix d'achat est compté pour un coût nul.}

	\fonc{Services}{Le boss enregistre ses services : nom, catégorie, prix de référence et durée indicative. Ils n'ont ni prix d'achat ni stock. Ils s'utilisent comme les produits lors de la saisie d'une opération.}

	\fonc{Catégories}{Produits et services sont classés par catégories créées et modifiées par le boss. Elles servent à la saisie rapide et aux statistiques.}

	\fonc{Stock optionnel}{Le module de stock est désactivé par défaut. Le boss peut l'activer globalement ou produit par produit, et l'arrêter à tout moment. Quand il active le suivi d'un produit, il saisit la quantité de départ, qui est enregistrée comme un premier comptage. Les quantités sont des nombres entiers, visibles par les employés, mais seul le boss peut les modifier.}

	\fonc{Mouvements de stock}{Le boss ne modifie jamais directement un chiffre de stock : il enregistre des \emph{mouvements}. Une entrée de marchandise (+20), une perte ou une casse ($-1$) s'ajoutent au stock. Un \emph{comptage} (« à 10~h, il y a 10 articles sur l'étagère ») fixe la quantité à un instant donné. Les ventes et leurs corrections génèrent elles-mêmes leurs mouvements. Le stock est calculé à partir de ces mouvements (voir la section~\ref{sec:stock}), et le boss peut toujours consulter la liste des mouvements qui expliquent la quantité d'un produit.}

	\fonc{Stock flexible}{La vente n'est jamais bloquée. Si le stock tombe sous le seuil d'alerte ou à zéro, une alerte s'affiche au boss uniquement (les employés ne reçoivent aucune alerte de stock) et le stock peut devenir négatif. Les mouvements saisis par le boss sont tracés dans l'audit.}

	\subsection{Enregistrement de l'activité}

	\fonc{Saisie rapide}{Le boss et les employés enregistrent une vente ou un service en choisissant l'article, la quantité et le prix. L'écran est pensé pour une saisie en quelques secondes, avec les articles fréquents en accès direct.}

	\fonc{Prix négociable}{Le prix de référence est pré-rempli mais modifiable librement, à la hausse comme à la baisse. L'écart (remise ou majoration) est conservé pour les statistiques.}

	\fonc{Panier multi-lignes}{Une même opération peut contenir plusieurs produits et services. Le total est calculé automatiquement et peut être ajusté globalement.}

	\fonc{Client facultatif}{Par défaut, l'opération est enregistrée pour un client anonyme. On peut rattacher un client (nom, téléphone), y compris après coup, mais ce n'est jamais obligatoire.}

	\fonc{Mode de paiement et note}{Un mode de paiement (espèces, mobile money, carte\dots) peut être indiqué, ainsi qu'une note libre. Une opération a un seul mode de paiement. La liste des modes de paiement est définie par le boss.}

	\fonc{Vente saisie par le boss pour un employé}{Le formulaire du boss comporte un champ « employé ayant fait la vente ». S'il est vide, la vente est considérée comme réalisée par le boss lui-même. L'opération est créée directement chez le boss, donc déjà synchronisée : seul le boss peut la modifier. L'historique et l'audit indiquent clairement qu'elle a été saisie par le boss pour le compte de l'employé. Elle est envoyée à l'employé uniquement en mode conservation.}

	\fonc{Correction et annulation}{Une opération est annulée avec un motif ou corrigée, sans effacement silencieux. Un employé corrige ses opérations tant qu'elles ne sont pas synchronisées. Après synchronisation, il ne peut plus rien modifier, quel que soit le mode de conservation : seul le boss corrige ou annule, avec un motif tracé dans l'audit. Il peut modifier toutes les opérations, les siennes comme celles des employés.}

	\fonc{Historique personnel}{Avant synchronisation, chaque utilisateur voit toutes ses opérations en attente. Après synchronisation, ce qu'un employé voit dépend du réglage de conservation (voir plus loin) ; il ne voit en aucun cas les opérations des autres. Le boss voit celles de toute l'équipe, filtrables par date, article, client ou employé.}

	\fonc{Récapitulatif de l'employé}{L'employé voit ses ventes du jour, de la semaine et du mois à partir des opérations qu'il peut voir. Si l'historique n'est pas conservé, seules les opérations en attente de synchronisation y figurent. Il n'accède pas aux chiffres des autres ni au tableau de bord global.}

	\subsection{Mode rafale}

	Le mode rafale sert à enregistrer très vite de nombreuses ventes, par exemple lorsqu'un coiffeur enchaîne les clients.

	\fonc{Modèle de rafale}{Avant de commencer, l'utilisateur définit un modèle : un ou plusieurs articles (produits, services, ou les deux), leurs quantités, leurs prix de référence et un mode de paiement. Il peut le modifier à tout moment.}

	\fonc{Une vente par clic}{Chaque clic sur le bouton enregistre immédiatement une vente complète contenant tous les articles du modèle, sans reremplir de formulaire. Chaque vente est indépendante des autres et concerne un client, qui reste facultatif.}

	\fonc{Souplesse}{Chaque vente d'une rafale peut être personnalisée ensuite dans l'historique (prix, articles, client, note), avec les mêmes règles que les autres opérations. Un bouton « annuler la dernière vente » évite les erreurs de double clic. Chaque vente garde l'heure de son clic.}

	\fonc{Disponibilité}{Le mode rafale est offert au boss et aux employés. Le boss peut y utiliser le champ « employé ayant fait la vente ».}

	\subsection{Conservation de l'historique côté employé}

	\fonc{Réglage du boss}{Le boss choisit si les opérations enregistrées par les employés disparaissent de leur téléphone après synchronisation, ou si elles sont conservées. Le réglage est global : il s'applique à tous les employés.}

	\fonc{Mode conservation}{Sur chaque appareil, l'employé retrouve les opérations saisies sur cet appareil depuis son appairage, ainsi que celles saisies par le boss à son nom, rangées par exercice (l'exercice attribué par le boss à la réception). Les opérations synchronisées et les exercices clos sont en lecture seule. Les corrections et annulations faites par le boss sur ses opérations, ainsi que les opérations saisies par le boss à son nom, lui sont renvoyées. L'employé ne voit jamais les dépenses, le tableau de bord, ni les opérations des autres.}

	\fonc{Mode disparition}{Après synchronisation, \textbf{tout} ce qui concerne l'opération disparaît du téléphone de l'employé : il ne voit plus rien des ventes envoyées, ni dans l'historique, ni dans le récapitulatif. Les opérations saisies par le boss pour lui ne lui sont pas transmises. Seuls subsistent le journal de ses synchronisations (date, nombre d'éléments, résultat) et la copie de secours chiffrée, qu'il ne peut ni lire ni voir.}

	\fonc{Prise d'effet}{Quand le boss change le réglage, il s'applique à la synchronisation suivante de chaque employé. Passer à la disparition efface tout l'historique existant sur son téléphone. Repasser à la conservation ne restitue rien : seules les nouvelles données sont conservées.}

	\fonc{Sécurité des données}{Une opération est considérée comme synchronisée uniquement après l'accusé de réception du téléphone du boss. Elle ne disparaît du téléphone de l'employé qu'à ce moment, afin qu'une coupure ou un passage par le serveur relais ne fasse rien perdre.}

	\subsection{Dépenses}

	\fonc{Dépenses supplémentaires}{Le boss enregistre les dépenses de l'entreprise (loyer, facture d'électricité, chaise cassée\dots) par un formulaire simple : nom, description, montant et date (aujourd'hui par défaut). Une dépense est rattachée à l'exercice ouvert au moment de sa saisie ; sa date est informative.}

	\fonc{Gestion}{Une dépense peut être modifiée ou annulée par le boss, avec trace dans l'audit. Les dépenses sont réservées au boss, saisies uniquement sur son téléphone et jamais transmises aux employés lors des synchronisations. Elles sont déduites dans le calcul du bénéfice.}

	\subsection{Synchronisation}

	\fonc{Toujours à la demande}{Aucune synchronisation n'est automatique. C'est toujours l'utilisateur qui la lance, application ouverte et au premier plan, et il reste sur l'écran de synchronisation jusqu'à la fin. Rien ne se passe en arrière-plan.}

	\fonc{Synchronisation sur réseau local (mode principal)}{Le boss ouvre l'écran « Recevoir une synchronisation » : son téléphone affiche un QR code contenant son adresse sur le réseau et un jeton de session à usage unique. L'employé scanne ce QR code et l'échange démarre, chiffré, directement entre les deux téléphones. Il suffit qu'ils soient sur le même réseau : le Wi-Fi de l'entreprise ou, à défaut, le partage de connexion activé sur le téléphone du boss. Internet n'est pas nécessaire. Aucune synchronisation n'est acceptée si aucun exercice n'est ouvert chez le boss.}

	\fonc{Sens de transfert}{Des employés vers le boss : opérations, clients créés, modifications de profil. Du boss vers les employés : catalogue (sans les prix d'achat, avec les quantités en stock), clients, le compte du propriétaire de l'appareil (dont ses réinitialisations de mot de passe), paramètres (dont le réglage de conservation), apparence de l'entreprise, nom et état des exercices, accusés de réception et, en mode conservation uniquement, les opérations saisies à leur nom et les corrections faites sur leurs opérations. Les dépenses, les prix d'achat, les mouvements de stock et le contenu des exercices clos ne sont jamais envoyés.}

	\fonc{Serveur relais (complément)}{Quand le boss est absent, un employé peut, s'il a internet, déposer ses opérations sur un petit serveur relais. Le relais sert de boîte aux lettres : il ne conserve que des données chiffrées de bout en bout, qu'il ne peut pas lire. Le boss récupère ces dépôts quand il lance lui-même une synchronisation par le relais, et y dépose en retour ses mises à jour et ses accusés de réception, que l'employé récupère à son tour. Une opération déposée sur le relais n'est considérée comme synchronisée qu'à la réception de l'accusé du boss. L'application fonctionne entièrement sans relais.}

	\fonc{Incrémentale et fiable}{Seuls les changements depuis la dernière synchronisation sont échangés. Une coupure en cours de route ne perd ni ne duplique aucune donnée : ce qui n'a pas été confirmé est renvoyé la fois suivante, et ce qui a déjà été reçu est ignoré.}

	\fonc{Suivi de synchronisation}{Chaque appareil affiche la date de dernière synchronisation, le nombre d'éléments en attente et le résultat. Chaque utilisateur consulte ainsi l'historique de ses propres synchronisations, et le boss celui de tous les employés.}

	\subsection{Tableau de bord du boss}

	Le tableau de bord porte sur l'exercice courant et se présente sous forme de totaux et de listes.

	\fonc{Chiffre d'affaires et bénéfice}{Chiffre d'affaires de l'exercice, du jour, de la semaine, du mois ou d'une période choisie, réparti entre produits et services, avec comparaison à la période précédente dans le même exercice. Le bénéfice indicatif est le chiffre d'affaires moins le coût des articles vendus et moins les dépenses.}

	\fonc{Coût des articles vendus}{Le coût est compté à la vente : prix d'achat $\times$ quantité vendue. Comme l'employé ne connaît pas le prix d'achat, c'est le téléphone du boss qui le fixe à la réception de l'opération, avec le prix d'achat en vigueur à ce moment-là.}

	\fonc{Dépenses}{Total des dépenses de l'exercice et liste détaillée.}

	\fonc{Activité des employés}{Pour chaque employé : nombre d'opérations, chiffre généré et part qu'il représente dans le chiffre d'affaires de l'entreprise (en pourcentage). Le boss compte comme un producteur.}

	\fonc{Produits et services}{Classement des articles les plus vendus et de ceux qui ne se vendent pas. Détail par catégorie.}

	\fonc{Remises et prix négociés}{Total des remises accordées, par employé et par article. Cela montre l'écart entre prix de référence et prix réellement pratiqués.}

	\fonc{Alertes}{Stock bas ou négatif, employé n'ayant pas synchronisé depuis un certain nombre de jours, opérations annulées ou corrigées récemment.}

	\fonc{Bilan simple}{Récapitulatif de l'activité d'un exercice ou d'une période, exportable. Il est donné à titre indicatif et n'a pas de valeur comptable.}

	\subsection{Sauvegarde, copie de secours et restauration}

	\fonc{Sauvegarde chiffrée}{Le boss sauvegarde ses données, exercices archivés inclus, vers son Google Drive ou iCloud, ou dans un fichier local qu'il peut conserver sur son ordinateur. La sauvegarde est chiffrée et peut être automatique. Elle contient aussi la clé privée du boss, protégée par sa clé de récupération.}

	\fonc{Copie de secours chez les employés}{Le téléphone du boss est le centre de l'entreprise : s'il est perdu, ce qu'il a reçu depuis sa dernière sauvegarde le serait aussi. Pour l'éviter, chaque opération envoyée par un employé est aussi gardée sur son téléphone, chiffrée avec la clé publique du boss, pendant un nombre de jours réglé par le boss (30 jours par défaut). L'employé ne peut ni la lire, ni la voir, ni l'exporter : seule la clé privée du boss permet de la déchiffrer. Elle est effacée automatiquement à l'expiration du délai, et immédiatement en cas de révocation de l'appareil.}

	\fonc{Restauration}{En cas de perte du téléphone, le boss restaure l'application sur un nouvel appareil à partir d'une sauvegarde, avec sa clé de récupération. À leur synchronisation suivante, les téléphones des employés renvoient leur copie de secours : le boss la déchiffre et récupère les opérations reçues après sa sauvegarde, sans doublon. Ces opérations rejoignent l'exercice ouvert, comme toute opération reçue. Les opérations saisies directement sur le téléphone du boss depuis sa dernière sauvegarde ne peuvent pas être récupérées, d'où l'intérêt de la sauvegarde automatique.}

	\fonc{Pas de sauvegarde sur le serveur}{Le serveur relais reste le plus léger possible : il ne stocke jamais de sauvegarde et ne conserve les messages que le temps nécessaire à leur récupération.}

	\subsection{Personnalisation et paramètres}

	\fonc{Apparence de l'entreprise}{Logo, couleurs et nom définis par le boss. Ces éléments sont diffusés aux employés à la synchronisation.}

	\fonc{Langues}{L'application est disponible en français et en anglais. Chaque utilisateur choisit sa langue ; d'autres langues pourront s'ajouter.}

	\fonc{Journal d'audit}{Le boss peut consulter qui a fait quoi et quand : opérations (en distinguant l'auteur de la saisie et l'employé concerné), corrections, annulations, dépenses, mouvements de stock, exercices, changements de catalogue et de réglages. L'audit est global et couvre tous les exercices.}

	%-------------------------------------------------
	\section{Règles de synchronisation}
	%-------------------------------------------------
	Le modèle est \textbf{en étoile} : le téléphone du boss est le centre. Un employé ne reçoit jamais les opérations de ses collègues, et deux employés ne se synchronisent pas entre eux. Les règles suivantes s'appliquent.

	\begin{itemize}[leftmargin=*]
		\item \textbf{Déclenchement.} Toute synchronisation est lancée manuellement par l'utilisateur, au premier plan. Le réseau local est le canal principal ; le relais est un complément facultatif.
		\item \textbf{Prix et catalogue.} Une vente garde le prix de vente et les informations en vigueur au moment où elle a été enregistrée, même si le boss modifie le catalogue ensuite. Le coût d'achat fait exception : il est fixé à la réception de l'opération par le boss.
		\item \textbf{Heures.} L'heure retenue pour une opération est celle de l'appareil de l'utilisateur. Elle sert à l'historique, aux statistiques par période et au calcul du stock, mais pas au choix de l'exercice. Aucune correction n'est exigée si les horloges divergent.
		\item \textbf{Exercice.} Une opération reçue par le boss rejoint l'exercice ouvert à ce moment. Un exercice clos ne reçoit plus rien. Sans exercice ouvert, aucune synchronisation n'est possible.
		\item \textbf{Propriété des données.} Les opérations appartiennent à l'employé auquel elles sont attribuées, même si le boss les a saisies à sa place ; le catalogue, les comptes, les dépenses, le stock et les paramètres appartiennent au boss, dont la version fait foi.
		\item \textbf{Opération synchronisée.} Une opération est synchronisée à l'accusé de réception du boss. Avant cet accusé, l'employé peut la corriger et la nouvelle version est envoyée. Après, il ne peut plus rien modifier, quel que soit le mode de conservation ; seul le boss peut la corriger ou l'annuler, avec un motif.
		\item \textbf{Retour au boss.} Les corrections du boss ne sont envoyées à l'employé qu'en mode conservation, et ne concernent que ses propres opérations : chaque appareil reçoit les corrections des opérations saisies sur lui, et les opérations saisies par le boss au nom de son propriétaire.
		\item \textbf{Un appareil, un utilisateur.} Chaque appareil appartient à un seul utilisateur et n'envoie que ses opérations. Les appareils d'un même utilisateur ne se synchronisent jamais entre eux.
		\item \textbf{Profil et mot de passe.} Ils appartiennent à l'utilisateur. Une modification passe par le boss avant d'atteindre ses autres appareils ; en cas de modifications sur deux appareils, la dernière reçue par le boss l'emporte.
		\item \textbf{Réglage de conservation.} Le mode appliqué est celui en vigueur au moment de la synchronisation.
		\item \textbf{Absence de doublons.} Chaque élément possède un identifiant unique et un numéro de version, ce qui évite les doublons en cas de nouvelle tentative ou de renvoi de la copie de secours.
		\item \textbf{Stock.} Le stock fait foi chez le boss et se calcule à partir des mouvements (section~\ref{sec:stock}). Les employés n'envoient jamais de quantité, seulement leurs ventes. En cas de ventes simultanées hors ligne, le stock peut devenir négatif ; une alerte est alors affichée au boss uniquement. Le stock n'est pas rattaché à un exercice.
		\item \textbf{Appareil révoqué.} Un appareil révoqué efface toutes ses données locales dès sa synchronisation suivante. Désactiver un compte révoque tous les appareils de ce compte.
	\end{itemize}

	%-------------------------------------------------
	\section{Calcul du stock}\label{sec:stock}
	%-------------------------------------------------
	Hors ligne, les informations n'arrivent pas chez le boss dans l'ordre où les événements se sont produits. Si le stock était un simple chiffre que l'on modifie, le résultat dépendrait de l'ordre d'arrivée. Exemple : à 9~h, un employé vend 2 articles hors ligne ; à 10~h, le boss compte 10 articles sur l'étagère et saisit 10 ; à 18~h, la vente de 9~h arrive et serait retirée une seconde fois, donnant 8 au lieu de 10. Link calcule donc le stock à partir de \textbf{mouvements datés}, conservés sur le téléphone du boss uniquement.

	\begin{itemize}[leftmargin=*]
		\item \textbf{Mouvements qui s'additionnent} : vente ($-$quantité vendue), correction ou annulation d'une vente (mouvement inverse ou différence), entrée de marchandise, perte ou casse. Une addition donne le même résultat quel que soit l'ordre d'arrivée.
		\item \textbf{Comptage} : « à telle heure, il y avait telle quantité ». C'est le seul mouvement qui fixe une valeur, et il est daté. La quantité de départ, à l'activation du stock d'un produit, est un comptage.
		\item \textbf{Règle de calcul} : stock d'un produit = quantité du dernier comptage + somme des mouvements dont l'heure est postérieure à ce comptage. Dans l'exemple, la vente de 9~h est antérieure au comptage de 10~h : elle y est déjà incluse et n'est pas retirée à nouveau. Le stock vaut 10, quel que soit le moment de la synchronisation.
		\item \textbf{Valeur affichée chez le boss} : la quantité calculée est gardée en mémoire dans la fiche du produit et mise à jour à chaque nouveau mouvement. Elle peut toujours être recalculée à partir des mouvements.
		\item \textbf{Valeur affichée chez l'employé} : quantité reçue du boss à la dernière synchronisation de cet appareil, moins les quantités des ventes saisies sur cet appareil et pas encore confirmées par le boss. Dès qu'une vente est confirmée, elle est incluse dans la quantité du boss et ne compte plus dans les ventes en attente : elle n'est ni comptée deux fois, ni oubliée.
		\item \textbf{Limite} : la règle s'appuie sur l'heure des appareils. Une vente faite très près d'un comptage, sur un téléphone dont l'horloge est décalée, peut tomber du mauvais côté. Au moment de saisir un comptage, l'application conseille donc au boss de synchroniser d'abord tous les employés.
	\end{itemize}

	%-------------------------------------------------
	\section{Sécurité et confidentialité}
	%-------------------------------------------------
	\begin{itemize}[leftmargin=*]
		\item Base de données locale chiffrée sur chaque téléphone.
		\item Mots de passe stockés de manière non réversible, avec un algorithme lent (Argon2id), car ils sont présents sur les appareils pour permettre la connexion hors ligne ; mot de passe provisoire à usage unique.
		\item Un appareil n'est lié qu'à un seul compte et ne contient que ce compte : aucun appareil ne détient les informations de connexion d'un autre employé. Le compte boss n'est utilisable que sur le téléphone du boss.
		\item Clés de chiffrement échangées lors de l'appairage par QR code ; chaque session de synchronisation locale est protégée par un jeton à usage unique.
		\item Échanges de synchronisation chiffrés, y compris via le serveur relais, qui ne peut pas lire les données.
		\item Copie de secours des employés chiffrée avec la clé publique du boss : illisible pour l'employé.
		\item Appairage contrôlé par le boss ; révocation possible à tout moment.
		\item Droits stricts : un employé ne voit que ses propres opérations, et ni la gestion des comptes, ni le tableau de bord global, ni les dépenses, ni les prix d'achat, ni les données de ses collègues.
		\item Dépenses, prix d'achat, mouvements de stock et contenu des exercices clos ne quittent jamais le téléphone du boss (hors sauvegarde chiffrée).
		\item Sauvegardes chiffrées, protégées par une clé de récupération détenue par le boss.
	\end{itemize}

	%-------------------------------------------------
	\section{Contraintes et choix techniques}
	%-------------------------------------------------
	\begin{itemize}[leftmargin=*]
		\item \textbf{Plateformes :} mobile uniquement, iOS et Android, sur les téléphones personnels des utilisateurs.
		\item \textbf{Technologie :} application React Native, avec une base de données locale SQLite chiffrée. L'hébergement du serveur relais est à définir plus tard.
		\item \textbf{Hors ligne d'abord :} toutes les fonctions marchent sans internet. La synchronisation se fait sur le réseau local (Wi-Fi ou partage de connexion) ; internet n'est utile que pour le relais, facultatif, et les sauvegardes vers le Drive.
		\item \textbf{Montants et quantités :} francs CFA, stockés en nombres entiers, sans virgule. Les quantités sont aussi des entiers.
		\item \textbf{Envoi des employés vers le boss :} chaque changement est inscrit dans un journal local, en ajout seul. Seuls les éléments explicitement confirmés par le boss reçoivent une date d'accusé de réception ; tout le reste est renvoyé à la synchronisation suivante. Aucun état « en cours » n'est conservé en base : si l'application s'interrompt, l'élément est simplement renvoyé, et le boss ignore ce qu'il a déjà reçu (même identifiant, même version). Après accusé, l'élément reste dans le journal, chiffré, le temps de la copie de secours.
		\item \textbf{Envoi du boss vers les employés :} chaque modification d'une donnée partagée reçoit chez le boss un numéro croissant (compteur de modification). Chaque appareil d'employé retient le dernier numéro reçu ; à la synchronisation, il l'annonce et reçoit tout ce qui a changé depuis, filtré selon ses droits. Ce numéro est enregistré dans la même transaction que les données reçues. Un nouvel appareil reçoit une copie complète. Si un employé annonce un numéro supérieur à celui du boss (boss restauré depuis une sauvegarde), il reçoit aussi une copie complète.
		\item \textbf{Données :} identifiants uniques (UUID) et horodatage pour fiabiliser la synchronisation. Chaque vente et dépense porte, chez le boss, l'identifiant de son exercice.
		\item \textbf{Évolutivité :} chaque donnée est rattachée à un identifiant d'entreprise, pour permettre plus tard plusieurs entreprises sans refonte.
		\item \textbf{Serveur relais :} facultatif, le plus léger possible, sans stockage durable ni accès aux données en clair.
		\item \textbf{Langues :} textes externalisés pour faciliter l'ajout de langues.
		\item \textbf{Responsabilité de l'utilisateur :} chacun doit garder son application à jour par les synchronisations et les mises à jour des stores.
	\end{itemize}

	%-------------------------------------------------
	\section{Éléments mis de côté pour l'instant}\label{sec:miscote}
	%-------------------------------------------------
	\subsection{Reportés à une version ultérieure}
	Ces éléments ont été identifiés pendant la conception mais ne font pas partie de cette version.
	\begin{itemize}[leftmargin=*]
		\item \textbf{Fusion des clients en double.} Deux employés hors ligne peuvent créer le même client ; aucun outil de fusion n'est prévu pour l'instant.
		\item \textbf{Autres devises et montants à virgule.} Seul le franc CFA est géré, en nombres entiers. Le format de stockage permettra d'ajouter d'autres devises (en stockant la plus petite unité, par exemple les centimes).
		\item \textbf{Quantités décimales.} Vente au poids ou au mètre : les quantités sont entières pour l'instant.
		\item \textbf{Paiement fractionné.} Une opération n'a qu'un seul mode de paiement.
	\end{itemize}

	\subsection{Écartés volontairement}
	Ces options ont été étudiées et rejetées au profit d'un autre choix.
	\begin{itemize}[leftmargin=*]
		\item \textbf{Synchronisation automatique ou en arrière-plan}, y compris à la connexion au réseau de l'entreprise : la synchronisation est toujours lancée par l'utilisateur.
		\item \textbf{Découverte automatique des appareils sur le réseau local} : remplacée par le QR code de session affiché par le boss.
		\item \textbf{Serveur relais comme canal principal} : Link reste pensé d'abord pour le hors ligne ; le relais n'est qu'un complément.
		\item \textbf{Rattachement d'une opération à l'exercice de sa date} (avec opérations « tardives » ajoutées à un exercice clos) : remplacé par le rattachement à l'exercice ouvert à la réception, qui garde les exercices clos strictement en lecture seule.
		\item \textbf{Stock modifié directement par une quantité} : remplacé par le calcul à partir des mouvements datés.
		\item \textbf{Téléphone partagé entre plusieurs employés} : chaque appareil appartient à un seul utilisateur. Ce cas est peu courant et compliquait la connexion hors ligne, la désactivation des comptes et la confidentialité entre collègues.
		\item \textbf{Connexion du boss sur un autre appareil} : le compte boss reste lié à son propre téléphone, centre de l'entreprise.
	\end{itemize}

	%-------------------------------------------------
	\section{Évolutions futures}
	%-------------------------------------------------
	\begin{itemize}[leftmargin=*]
		\item Graphiques : évolution du chiffre d'affaires, heures de pointe, répartition par mode de paiement.
		\item Gestion de plusieurs entreprises par un même boss, et employés rattachés à plusieurs entreprises (déjà anticipé dans le modèle de données).
		\item Rôle intermédiaire (adjoint) avec accès partiel au tableau de bord.
		\item Historique des prix d'achat datés, pour un coût plus exact.
		\item Restitution de l'historique à un employé après changement de téléphone.
		\item Comparaison entre exercices.
		\item Autres langues et personnalisations supplémentaires.
		\item Export plus complet des bilans et statistiques.
	\end{itemize}

\end{document}
